% --- Archon physics formalization source begin ---
% archon:physics
% archon:covers PhyXMiniProblems/problem_phyx_mini_0920.lean
% archon:source-report reports/phyx_mini/problem_phyx_mini_0920.source.json

\chapter{Physics problem phyx\_mini\_0920}
\label{ch:PhyXMiniProblems_problem_phyx_mini_0920}

\paragraph{Problem source.}
\#\# Physical scenario

Figure shows a conducting disk with radius \textbackslash{}( R \textbackslash{}) that lies in the \textbackslash{}( xy \textbackslash{})-plane and rotates with constant angular velocity \textbackslash{}( \textbackslash{}omega \textbackslash{}) about the \textbackslash{}( z \textbackslash{})-axis. The disk is in a uniform, constant \textbackslash{}( \textbackslash{}vec\{B\} \textbackslash{}) field in the \textbackslash{}( z \textbackslash{})-direction.

\#\# Question

Find the induced emf between the center and the rim of the disk.

\#\# Answer choices

- A: A: \textbackslash{}frac\{1\}\{4\} \textbackslash{}omega BR\textasciicircum{}2

- B: B: \textbackslash{}frac\{1\}\{2\} \textbackslash{}omega BR

- C: C: \textbackslash{}frac\{1\}\{2\} \textbackslash{}omega BR\textasciicircum{}2

- D: D: \textbackslash{}omega BR\textasciicircum{}2

\#\# Figure caption (auxiliary; use the image as primary evidence)

The diagram illustrates a rotating circular disk with the following elements:

1. **Disk**: The disk is rotating around its central axis, labeled with an angular velocity symbol (\textbackslash{}(\textbackslash{}omega\textbackslash{})) in the \textbackslash{}(z\textbackslash{})-axis direction. The disk is positioned perpendicular to this axis.

2. **Axes**: The \textbackslash{}(x\textbackslash{}), \textbackslash{}(y\textbackslash{}), and \textbackslash{}(z\textbackslash{}) axes are shown for reference. The \textbackslash{}(z\textbackslash{})-axis is vertical, perpendicular to the plane of the disk, while the \textbackslash{}(x\textbackslash{}) and \textbackslash{}(y\textbackslash{}) axes lie in the plane of the disk.

3. **Magnetic Field (\textbackslash{}(\textbackslash{}vec\{B\}\textbackslash{}))**: Several blue arrows labeled \textbackslash{}(\textbackslash{}vec\{B\}\textbackslash{}) represent a magnetic field pointing into the surface.

4. **Radial Segment**: A small radial segment of the disk is highlighted. This segment extends a distance \textbackslash{}(dr\textbackslash{}) from the center of the disk to a point at radius \textbackslash{}(r\textbackslash{}). It is shown with a green arrow labeled \textbackslash{}(\textbackslash{}vec\{v\}\textbackslash{}) indicating its tangential velocity, which is \textbackslash{}(\textbackslash{}nu = \textbackslash{}omega r\textbackslash{}).

5. **Text Descriptions**:
   - "Speed of small radial segment of length \textbackslash{}(dr\textbackslash{}) at radius \textbackslash{}(r\textbackslash{}) is \textbackslash{}(\textbackslash{}nu = \textbackslash{}omega r\textbackslash{})."
   - "Emf induced across this segment is \textbackslash{}(d\textbackslash{}mathcal\{E\} = vB \textbackslash{}, dr = \textbackslash{}omega Br \textbackslash{}, dr\textbackslash{})." This indicates the electromotive force induced in the segment due to its motion in the magnetic field.

6. **Connections**:
   - The disk is connected via conducting brushes indicated by \textbackslash{}(b\textbackslash{}), showing connections to an external circuit with a resistor symbol and a purple arrow indicating current (\textbackslash{}(I\textbackslash{})) flow. 

The diagram explains the physical scenario of electromagnetic induction in a rotating disk within a magnetic field, analyzing the induced electromotive force and current flow.

\#\# Dataset metadata

Electromagnetic Induction; Physical Model Grounding Reasoning, Multi-Formula Reasoning

\paragraph{Recorded answer/context.}
C: C: \textbackslash{}frac\{1\}\{2\} \textbackslash{}omega BR\textasciicircum{}2

\paragraph{Figure/image path.}
/root/proposal\_for\_physic/science-mango/phyx\_mini\_run/phyx\_data/test\_image/920.png

\paragraph{Formalization target.}
create a compiling Lean file with sorry bodies at `PhyXMiniProblems/problem\_phyx\_mini\_0920.lean`. The Lean declarations must preserve the physical quantities, dimensions or dimensional roles, figure labels, governing-law hypotheses, and final relation expressed by this problem.
Use Mathlib/Physlib names found through LeanExplore where available. If a physics API is missing, introduce faithful local abstractions rather than scalar placeholder aliases.

\begin{theorem}[Physics formalization target]
\label{thm:physics:phyx_mini_0920:target}
\lean{PhyXMiniProblems.ProblemPhyXMini0920.problem_phyx_mini_0920}
\uses{def:physics:phyx-mini-0920:phyxminiproblems-problemphyxmini0920-lengthinmeters, def:physics:phyx-mini-0920:phyxminiproblems-problemphyxmini0920-angularspeedinradianspersecond, def:physics:phyx-mini-0920:phyxminiproblems-problemphyxmini0920-magneticfluxdensityinteslas, def:physics:phyx-mini-0920:phyxminiproblems-problemphyxmini0920-emfinvolts, def:physics:phyx-mini-0920:phyxminiproblems-problemphyxmini0920-rotatingconductingdisksetup, def:physics:phyx-mini-0920:phyxminiproblems-problemphyxmini0920-matchesrotatingconductingdiskdescription, def:physics:phyx-mini-0920:phyxminiproblems-problemphyxmini0920-matchessuppliedrotatingdiskfigure, def:physics:phyx-mini-0920:phyxminiproblems-problemphyxmini0920-hasphysicalrotatingdiskparameters, def:physics:phyx-mini-0920:phyxminiproblems-problemphyxmini0920-satisfiesuniformaxialmagneticfield, def:physics:phyx-mini-0920:phyxminiproblems-problemphyxmini0920-satisfiesrotatingdiskmotionalemflaws}
The assigned autoformalize agent should translate this physics problem into Lean declarations in the covered file, with theorem and lemma proof bodies written as `by sorry`.
\end{theorem}
\begin{proof}
This is an autoformalization task, not a proof task. Produce statements that can later be proved by the physics prover without weakening the physical model.
\end{proof}

% --- Archon declaration topology begin ---
\section{Declaration topology}
\begin{definition}[angular Speed Dimension]
  \label{def:physics:phyx-mini-0920:phyxminiproblems-problemphyxmini0920-angularspeeddimension}
  \lean{PhyXMiniProblems.ProblemPhyXMini0920.angularSpeedDimension}
  Angular speed has dimension T⁻¹; radians are dimensionless
\end{definition}
\begin{proof}
  This is the declared model definition of angular Speed Dimension.
\end{proof}
\begin{definition}[speed Dimension]
  \label{def:physics:phyx-mini-0920:phyxminiproblems-problemphyxmini0920-speeddimension}
  \lean{PhyXMiniProblems.ProblemPhyXMini0920.speedDimension}
  Tangential speed has dimension L T⁻¹
\end{definition}
\begin{proof}
  This is the declared model definition of speed Dimension.
\end{proof}
\begin{definition}[magnetic Flux Density Dimension]
  \label{def:physics:phyx-mini-0920:phyxminiproblems-problemphyxmini0920-magneticfluxdensitydimension}
  \lean{PhyXMiniProblems.ProblemPhyXMini0920.magneticFluxDensityDimension}
  Magnetic flux density has dimension M T⁻¹ C⁻¹
\end{definition}
\begin{proof}
  This is the declared model definition of magnetic Flux Density Dimension.
\end{proof}
\begin{definition}[motional Emf Density Dimension]
  \label{def:physics:phyx-mini-0920:phyxminiproblems-problemphyxmini0920-motionalemfdensitydimension}
  \lean{PhyXMiniProblems.ProblemPhyXMini0920.motionalEmfDensityDimension}
  Motional-EMF density has electric-field dimension M L T⁻² C⁻¹
\end{definition}
\begin{proof}
  This is the declared model definition of motional Emf Density Dimension.
\end{proof}
\begin{definition}[emf Dimension]
  \label{def:physics:phyx-mini-0920:phyxminiproblems-problemphyxmini0920-emfdimension}
  \lean{PhyXMiniProblems.ProblemPhyXMini0920.emfDimension}
  EMF has electric-potential dimension M L² T⁻² C⁻¹
\end{definition}
\begin{proof}
  This is the declared model definition of emf Dimension.
\end{proof}
\begin{definition}[electric Current Dimension]
  \label{def:physics:phyx-mini-0920:phyxminiproblems-problemphyxmini0920-electriccurrentdimension}
  \lean{PhyXMiniProblems.ProblemPhyXMini0920.electricCurrentDimension}
  Electric current has dimension C T⁻¹
\end{definition}
\begin{proof}
  This is the declared model definition of electric Current Dimension.
\end{proof}
\begin{definition}[electrical Resistance Dimension]
  \label{def:physics:phyx-mini-0920:phyxminiproblems-problemphyxmini0920-electricalresistancedimension}
  \lean{PhyXMiniProblems.ProblemPhyXMini0920.electricalResistanceDimension}
  Electrical resistance has dimension M L² T⁻¹ C⁻²
\end{definition}
\begin{proof}
  This is the declared model definition of electrical Resistance Dimension.
\end{proof}
\begin{definition}[Length Quantity]
  \label{def:physics:phyx-mini-0920:phyxminiproblems-problemphyxmini0920-lengthquantity}
  \lean{PhyXMiniProblems.ProblemPhyXMini0920.LengthQuantity}
  A nonnegative, unit-independent physical length
\end{definition}
\begin{proof}
  This is the declared model definition of Length Quantity.
\end{proof}
\begin{definition}[Angular Speed Magnitude]
  \label{def:physics:phyx-mini-0920:phyxminiproblems-problemphyxmini0920-angularspeedmagnitude}
  \lean{PhyXMiniProblems.ProblemPhyXMini0920.AngularSpeedMagnitude}
  \uses{def:physics:phyx-mini-0920:phyxminiproblems-problemphyxmini0920-angularspeeddimension}
  A nonnegative, unit-independent angular-speed magnitude
\end{definition}
\begin{proof}
  This is the declared model definition of Angular Speed Magnitude.
\end{proof}
\begin{definition}[Speed Magnitude]
  \label{def:physics:phyx-mini-0920:phyxminiproblems-problemphyxmini0920-speedmagnitude}
  \lean{PhyXMiniProblems.ProblemPhyXMini0920.SpeedMagnitude}
  \uses{def:physics:phyx-mini-0920:phyxminiproblems-problemphyxmini0920-speeddimension}
  A nonnegative, unit-independent tangential-speed magnitude
\end{definition}
\begin{proof}
  This is the declared model definition of Speed Magnitude.
\end{proof}
\begin{definition}[Magnetic Flux Density Magnitude]
  \label{def:physics:phyx-mini-0920:phyxminiproblems-problemphyxmini0920-magneticfluxdensitymagnitude}
  \lean{PhyXMiniProblems.ProblemPhyXMini0920.MagneticFluxDensityMagnitude}
  \uses{def:physics:phyx-mini-0920:phyxminiproblems-problemphyxmini0920-magneticfluxdensitydimension}
  A nonnegative, unit-independent magnetic-flux-density magnitude
\end{definition}
\begin{proof}
  This is the declared model definition of Magnetic Flux Density Magnitude.
\end{proof}
\begin{definition}[Motional Emf Density Magnitude]
  \label{def:physics:phyx-mini-0920:phyxminiproblems-problemphyxmini0920-motionalemfdensitymagnitude}
  \lean{PhyXMiniProblems.ProblemPhyXMini0920.MotionalEmfDensityMagnitude}
  \uses{def:physics:phyx-mini-0920:phyxminiproblems-problemphyxmini0920-motionalemfdensitydimension}
  A nonnegative local motional-EMF density, measured in volts per metre
\end{definition}
\begin{proof}
  This is the declared model definition of Motional Emf Density Magnitude.
\end{proof}
\begin{definition}[Emf Magnitude]
  \label{def:physics:phyx-mini-0920:phyxminiproblems-problemphyxmini0920-emfmagnitude}
  \lean{PhyXMiniProblems.ProblemPhyXMini0920.EmfMagnitude}
  \uses{def:physics:phyx-mini-0920:phyxminiproblems-problemphyxmini0920-emfdimension}
  A nonnegative induced-EMF magnitude
\end{definition}
\begin{proof}
  This is the declared model definition of Emf Magnitude.
\end{proof}
\begin{definition}[Electric Current Magnitude]
  \label{def:physics:phyx-mini-0920:phyxminiproblems-problemphyxmini0920-electriccurrentmagnitude}
  \lean{PhyXMiniProblems.ProblemPhyXMini0920.ElectricCurrentMagnitude}
  \uses{def:physics:phyx-mini-0920:phyxminiproblems-problemphyxmini0920-electriccurrentdimension}
  A nonnegative electric-current magnitude in the depicted external circuit
\end{definition}
\begin{proof}
  This is the declared model definition of Electric Current Magnitude.
\end{proof}
\begin{definition}[Resistance Magnitude]
  \label{def:physics:phyx-mini-0920:phyxminiproblems-problemphyxmini0920-resistancemagnitude}
  \lean{PhyXMiniProblems.ProblemPhyXMini0920.ResistanceMagnitude}
  \uses{def:physics:phyx-mini-0920:phyxminiproblems-problemphyxmini0920-electricalresistancedimension}
  A nonnegative electrical resistance for the depicted load
\end{definition}
\begin{proof}
  This is the declared model definition of Resistance Magnitude.
\end{proof}
\begin{definition}[nonnegative SIReadout]
  \label{def:physics:phyx-mini-0920:phyxminiproblems-problemphyxmini0920-nonnegativesireadout}
  \lean{PhyXMiniProblems.ProblemPhyXMini0920.nonnegativeSIReadout}
  Read a nonnegative dimensionful quantity in coherent SI units
\end{definition}
\begin{proof}
  This is the declared model definition of nonnegative SIReadout.
\end{proof}
\begin{definition}[length In Meters]
  \label{def:physics:phyx-mini-0920:phyxminiproblems-problemphyxmini0920-lengthinmeters}
  \lean{PhyXMiniProblems.ProblemPhyXMini0920.lengthInMeters}
  \uses{def:physics:phyx-mini-0920:phyxminiproblems-problemphyxmini0920-lengthquantity, def:physics:phyx-mini-0920:phyxminiproblems-problemphyxmini0920-nonnegativesireadout}
  Read a physical length in metres
\end{definition}
\begin{proof}
  This is the declared model definition of length In Meters.
\end{proof}
\begin{definition}[angular Speed In Radians Per Second]
  \label{def:physics:phyx-mini-0920:phyxminiproblems-problemphyxmini0920-angularspeedinradianspersecond}
  \lean{PhyXMiniProblems.ProblemPhyXMini0920.angularSpeedInRadiansPerSecond}
  \uses{def:physics:phyx-mini-0920:phyxminiproblems-problemphyxmini0920-angularspeedmagnitude, def:physics:phyx-mini-0920:phyxminiproblems-problemphyxmini0920-nonnegativesireadout}
  Read angular speed in radians per second
\end{definition}
\begin{proof}
  This is the declared model definition of angular Speed In Radians Per Second.
\end{proof}
\begin{definition}[speed In Meters Per Second]
  \label{def:physics:phyx-mini-0920:phyxminiproblems-problemphyxmini0920-speedinmeterspersecond}
  \lean{PhyXMiniProblems.ProblemPhyXMini0920.speedInMetersPerSecond}
  \uses{def:physics:phyx-mini-0920:phyxminiproblems-problemphyxmini0920-speedmagnitude, def:physics:phyx-mini-0920:phyxminiproblems-problemphyxmini0920-nonnegativesireadout}
  Read tangential speed in metres per second
\end{definition}
\begin{proof}
  This is the declared model definition of speed In Meters Per Second.
\end{proof}
\begin{definition}[magnetic Flux Density In Teslas]
  \label{def:physics:phyx-mini-0920:phyxminiproblems-problemphyxmini0920-magneticfluxdensityinteslas}
  \lean{PhyXMiniProblems.ProblemPhyXMini0920.magneticFluxDensityInTeslas}
  \uses{def:physics:phyx-mini-0920:phyxminiproblems-problemphyxmini0920-magneticfluxdensitymagnitude, def:physics:phyx-mini-0920:phyxminiproblems-problemphyxmini0920-nonnegativesireadout}
  Read magnetic flux density in teslas
\end{definition}
\begin{proof}
  This is the declared model definition of magnetic Flux Density In Teslas.
\end{proof}
\begin{definition}[motional Emf Density In Volts Per Meter]
  \label{def:physics:phyx-mini-0920:phyxminiproblems-problemphyxmini0920-motionalemfdensityinvoltspermeter}
  \lean{PhyXMiniProblems.ProblemPhyXMini0920.motionalEmfDensityInVoltsPerMeter}
  \uses{def:physics:phyx-mini-0920:phyxminiproblems-problemphyxmini0920-motionalemfdensitymagnitude, def:physics:phyx-mini-0920:phyxminiproblems-problemphyxmini0920-nonnegativesireadout}
  Read local motional-EMF density in volts per metre
\end{definition}
\begin{proof}
  This is the declared model definition of motional Emf Density In Volts Per Meter.
\end{proof}
\begin{definition}[emf In Volts]
  \label{def:physics:phyx-mini-0920:phyxminiproblems-problemphyxmini0920-emfinvolts}
  \lean{PhyXMiniProblems.ProblemPhyXMini0920.emfInVolts}
  \uses{def:physics:phyx-mini-0920:phyxminiproblems-problemphyxmini0920-emfmagnitude, def:physics:phyx-mini-0920:phyxminiproblems-problemphyxmini0920-nonnegativesireadout}
  Read an induced-EMF magnitude in volts
\end{definition}
\begin{proof}
  This is the declared model definition of emf In Volts.
\end{proof}
\begin{definition}[current In Amperes]
  \label{def:physics:phyx-mini-0920:phyxminiproblems-problemphyxmini0920-currentinamperes}
  \lean{PhyXMiniProblems.ProblemPhyXMini0920.currentInAmperes}
  \uses{def:physics:phyx-mini-0920:phyxminiproblems-problemphyxmini0920-electriccurrentmagnitude, def:physics:phyx-mini-0920:phyxminiproblems-problemphyxmini0920-nonnegativesireadout}
  Read current magnitude in amperes
\end{definition}
\begin{proof}
  This is the declared model definition of current In Amperes.
\end{proof}
\begin{definition}[resistance In Ohms]
  \label{def:physics:phyx-mini-0920:phyxminiproblems-problemphyxmini0920-resistanceinohms}
  \lean{PhyXMiniProblems.ProblemPhyXMini0920.resistanceInOhms}
  \uses{def:physics:phyx-mini-0920:phyxminiproblems-problemphyxmini0920-resistancemagnitude, def:physics:phyx-mini-0920:phyxminiproblems-problemphyxmini0920-nonnegativesireadout}
  Read electrical resistance in ohms
\end{definition}
\begin{proof}
  This is the declared model definition of resistance In Ohms.
\end{proof}
\begin{definition}[Coordinate Axis]
  \label{def:physics:phyx-mini-0920:phyxminiproblems-problemphyxmini0920-coordinateaxis}
  \lean{PhyXMiniProblems.ProblemPhyXMini0920.CoordinateAxis}
  The three Cartesian axes explicitly labelled in the supplied raster
\end{definition}
\begin{proof}
  This is the declared model definition of Coordinate Axis.
\end{proof}
\begin{definition}[Coordinate Plane]
  \label{def:physics:phyx-mini-0920:phyxminiproblems-problemphyxmini0920-coordinateplane}
  \lean{PhyXMiniProblems.ProblemPhyXMini0920.CoordinatePlane}
  Coordinate plane occupied by the idealized disk
\end{definition}
\begin{proof}
  This is the declared model definition of Coordinate Plane.
\end{proof}
\begin{definition}[Axial Direction]
  \label{def:physics:phyx-mini-0920:phyxminiproblems-problemphyxmini0920-axialdirection}
  \lean{PhyXMiniProblems.ProblemPhyXMini0920.AxialDirection}
  Axial direction shared by the displayed rotation and magnetic field
\end{definition}
\begin{proof}
  This is the declared model definition of Axial Direction.
\end{proof}
\begin{definition}[Disk Material Model]
  \label{def:physics:phyx-mini-0920:phyxminiproblems-problemphyxmini0920-diskmaterialmodel}
  \lean{PhyXMiniProblems.ProblemPhyXMini0920.DiskMaterialModel}
  Material idealization assigned to the rotating object
\end{definition}
\begin{proof}
  This is the declared model definition of Disk Material Model.
\end{proof}
\begin{definition}[Disk Contact]
  \label{def:physics:phyx-mini-0920:phyxminiproblems-problemphyxmini0920-diskcontact}
  \lean{PhyXMiniProblems.ProblemPhyXMini0920.DiskContact}
  The two electrical contacts between which the EMF is requested
\end{definition}
\begin{proof}
  This is the declared model definition of Disk Contact.
\end{proof}
\begin{definition}[External Current Direction]
  \label{def:physics:phyx-mini-0920:phyxminiproblems-problemphyxmini0920-externalcurrentdirection}
  \lean{PhyXMiniProblems.ProblemPhyXMini0920.ExternalCurrentDirection}
  Direction of conventional current shown in the external load branch
\end{definition}
\begin{proof}
  This is the declared model definition of External Current Direction.
\end{proof}
\begin{definition}[Figure Quantity Label]
  \label{def:physics:phyx-mini-0920:phyxminiproblems-problemphyxmini0920-figurequantitylabel}
  \lean{PhyXMiniProblems.ProblemPhyXMini0920.FigureQuantityLabel}
  Symbolic quantity labels visible in image 920.png
\end{definition}
\begin{proof}
  This is the declared model definition of Figure Quantity Label.
\end{proof}
\begin{definition}[expected Printed Symbol]
  \label{def:physics:phyx-mini-0920:phyxminiproblems-problemphyxmini0920-expectedprintedsymbol}
  \lean{PhyXMiniProblems.ProblemPhyXMini0920.expectedPrintedSymbol}
  \uses{def:physics:phyx-mini-0920:phyxminiproblems-problemphyxmini0920-figurequantitylabel}
  Printed symbol expected for each labelled quantity in the raster
\end{definition}
\begin{proof}
  This is the declared model definition of expected Printed Symbol.
\end{proof}
\begin{definition}[Rotating Disk Figure]
  \label{def:physics:phyx-mini-0920:phyxminiproblems-problemphyxmini0920-rotatingdiskfigure}
  \lean{PhyXMiniProblems.ProblemPhyXMini0920.RotatingDiskFigure}
  \uses{def:physics:phyx-mini-0920:phyxminiproblems-problemphyxmini0920-coordinateaxis, def:physics:phyx-mini-0920:phyxminiproblems-problemphyxmini0920-diskcontact, def:physics:phyx-mini-0920:phyxminiproblems-problemphyxmini0920-externalcurrentdirection, def:physics:phyx-mini-0920:phyxminiproblems-problemphyxmini0920-figurequantitylabel}
  Literal presentation facts transcribed from the primary raster
\end{definition}
\begin{proof}
  This is the declared model definition of Rotating Disk Figure.
\end{proof}
\begin{definition}[Rotating Conducting Disk Setup]
  \label{def:physics:phyx-mini-0920:phyxminiproblems-problemphyxmini0920-rotatingconductingdisksetup}
  \lean{PhyXMiniProblems.ProblemPhyXMini0920.RotatingConductingDiskSetup}
  \uses{def:physics:phyx-mini-0920:phyxminiproblems-problemphyxmini0920-lengthquantity, def:physics:phyx-mini-0920:phyxminiproblems-problemphyxmini0920-angularspeedmagnitude, def:physics:phyx-mini-0920:phyxminiproblems-problemphyxmini0920-speedmagnitude, def:physics:phyx-mini-0920:phyxminiproblems-problemphyxmini0920-magneticfluxdensitymagnitude, def:physics:phyx-mini-0920:phyxminiproblems-problemphyxmini0920-motionalemfdensitymagnitude, def:physics:phyx-mini-0920:phyxminiproblems-problemphyxmini0920-emfmagnitude, def:physics:phyx-mini-0920:phyxminiproblems-problemphyxmini0920-electriccurrentmagnitude, def:physics:phyx-mini-0920:phyxminiproblems-problemphyxmini0920-resistancemagnitude, def:physics:phyx-mini-0920:phyxminiproblems-problemphyxmini0920-coordinateaxis, def:physics:phyx-mini-0920:phyxminiproblems-problemphyxmini0920-coordinateplane, def:physics:phyx-mini-0920:phyxminiproblems-problemphyxmini0920-axialdirection, def:physics:phyx-mini-0920:phyxminiproblems-problemphyxmini0920-diskmaterialmodel, def:physics:phyx-mini-0920:phyxminiproblems-problemphyxmini0920-diskcontact, def:physics:phyx-mini-0920:phyxminiproblems-problemphyxmini0920-rotatingdiskfigure}
  Independent physical quantities and profiles of the rotating-disk apparatus. The induced EMF is an observable indexed by its contacts; it is not defined from an answer expression. The radius-indexed speed and EMF-density fields return dimensionful quantities, with the index explicitly interpreted in metres by the governing laws below
\end{definition}
\begin{proof}
  This is the declared model definition of Rotating Conducting Disk Setup.
\end{proof}
\begin{definition}[z Axis Index]
  \label{def:physics:phyx-mini-0920:phyxminiproblems-problemphyxmini0920-zaxisindex}
  \lean{PhyXMiniProblems.ProblemPhyXMini0920.zAxisIndex}
  The standard z component in Physlib's three-dimensional field vectors
\end{definition}
\begin{proof}
  This is the declared model definition of z Axis Index.
\end{proof}
\begin{definition}[Matches Rotating Conducting Disk Description]
  \label{def:physics:phyx-mini-0920:phyxminiproblems-problemphyxmini0920-matchesrotatingconductingdiskdescription}
  \lean{PhyXMiniProblems.ProblemPhyXMini0920.MatchesRotatingConductingDiskDescription}
  \uses{def:physics:phyx-mini-0920:phyxminiproblems-problemphyxmini0920-rotatingconductingdisksetup}
  Qualitative assumptions stated in the problem prose
\end{definition}
\begin{proof}
  This is the declared model definition of Matches Rotating Conducting Disk Description.
\end{proof}
\begin{definition}[Matches Supplied Rotating Disk Figure]
  \label{def:physics:phyx-mini-0920:phyxminiproblems-problemphyxmini0920-matchessuppliedrotatingdiskfigure}
  \lean{PhyXMiniProblems.ProblemPhyXMini0920.MatchesSuppliedRotatingDiskFigure}
  \uses{def:physics:phyx-mini-0920:phyxminiproblems-problemphyxmini0920-expectedprintedsymbol, def:physics:phyx-mini-0920:phyxminiproblems-problemphyxmini0920-rotatingconductingdisksetup}
  Facts read from the primary image: all axes and symbolic labels, four axial field arrows, the local radial element and tangential velocity, and the two-brush external circuit. No EMF value or answer formula occurs here
\end{definition}
\begin{proof}
  This is the declared model definition of Matches Supplied Rotating Disk Figure.
\end{proof}
\begin{definition}[Has Physical Rotating Disk Parameters]
  \label{def:physics:phyx-mini-0920:phyxminiproblems-problemphyxmini0920-hasphysicalrotatingdiskparameters}
  \lean{PhyXMiniProblems.ProblemPhyXMini0920.HasPhysicalRotatingDiskParameters}
  \uses{def:physics:phyx-mini-0920:phyxminiproblems-problemphyxmini0920-lengthinmeters, def:physics:phyx-mini-0920:phyxminiproblems-problemphyxmini0920-angularspeedinradianspersecond, def:physics:phyx-mini-0920:phyxminiproblems-problemphyxmini0920-magneticfluxdensityinteslas, def:physics:phyx-mini-0920:phyxminiproblems-problemphyxmini0920-currentinamperes, def:physics:phyx-mini-0920:phyxminiproblems-problemphyxmini0920-resistanceinohms, def:physics:phyx-mini-0920:phyxminiproblems-problemphyxmini0920-rotatingconductingdisksetup}
  Positivity conditions selecting the nondegenerate physical apparatus
\end{definition}
\begin{proof}
  This is the declared model definition of Has Physical Rotating Disk Parameters.
\end{proof}
\begin{definition}[Satisfies Uniform Axial Magnetic Field]
  \label{def:physics:phyx-mini-0920:phyxminiproblems-problemphyxmini0920-satisfiesuniformaxialmagneticfield}
  \lean{PhyXMiniProblems.ProblemPhyXMini0920.SatisfiesUniformAxialMagneticField}
  \uses{def:physics:phyx-mini-0920:phyxminiproblems-problemphyxmini0920-magneticfluxdensityinteslas, def:physics:phyx-mini-0920:phyxminiproblems-problemphyxmini0920-rotatingconductingdisksetup, def:physics:phyx-mini-0920:phyxminiproblems-problemphyxmini0920-zaxisindex}
  The Physlib magnetic field is spatially uniform, time independent, and equal to B in its z component. This uses Physlib's spacetime-dependent Electromagnetism.MagneticField 3 rather than replacing the physical field by a bare scalar
\end{definition}
\begin{proof}
  This is the declared model definition of Satisfies Uniform Axial Magnetic Field.
\end{proof}
\begin{definition}[Satisfies Rotating Disk Motional Emf Laws]
  \label{def:physics:phyx-mini-0920:phyxminiproblems-problemphyxmini0920-satisfiesrotatingdiskmotionalemflaws}
  \lean{PhyXMiniProblems.ProblemPhyXMini0920.SatisfiesRotatingDiskMotionalEmfLaws}
  \uses{def:physics:phyx-mini-0920:phyxminiproblems-problemphyxmini0920-lengthinmeters, def:physics:phyx-mini-0920:phyxminiproblems-problemphyxmini0920-angularspeedinradianspersecond, def:physics:phyx-mini-0920:phyxminiproblems-problemphyxmini0920-speedinmeterspersecond, def:physics:phyx-mini-0920:phyxminiproblems-problemphyxmini0920-magneticfluxdensityinteslas, def:physics:phyx-mini-0920:phyxminiproblems-problemphyxmini0920-motionalemfdensityinvoltspermeter, def:physics:phyx-mini-0920:phyxminiproblems-problemphyxmini0920-emfinvolts, def:physics:phyx-mini-0920:phyxminiproblems-problemphyxmini0920-rotatingconductingdisksetup}
  Governing kinematic and motional-induction laws used by the displayed derivation. Rigid rotation gives the local tangential speed v(r) = ω r. With v perpendicular to the axial field, the local motional-EMF density is v(r) B, exactly the coefficient of dr in the raster's dℰ = v B dr. The requested center-to-rim EMF is the radial line integral of that density. These are local and integral laws. They contain no evaluated R² / 2 and no displayed answer choice
\end{definition}
\begin{proof}
  This is the declared model definition of Satisfies Rotating Disk Motional Emf Laws.
\end{proof}
\begin{definition}[Answer Choice]
  \label{def:physics:phyx-mini-0920:phyxminiproblems-problemphyxmini0920-answerchoice}
  \lean{PhyXMiniProblems.ProblemPhyXMini0920.AnswerChoice}
  Labels of the four answer choices printed with the problem
\end{definition}
\begin{proof}
  This is the declared model definition of Answer Choice.
\end{proof}
\begin{definition}[displayed Expression In SI]
  \label{def:physics:phyx-mini-0920:phyxminiproblems-problemphyxmini0920-answerchoice-displayedexpressioninsi}
  \lean{PhyXMiniProblems.ProblemPhyXMini0920.AnswerChoice.displayedExpressionInSI}
  \uses{def:physics:phyx-mini-0920:phyxminiproblems-problemphyxmini0920-lengthinmeters, def:physics:phyx-mini-0920:phyxminiproblems-problemphyxmini0920-angularspeedinradianspersecond, def:physics:phyx-mini-0920:phyxminiproblems-problemphyxmini0920-magneticfluxdensityinteslas, def:physics:phyx-mini-0920:phyxminiproblems-problemphyxmini0920-rotatingconductingdisksetup, def:physics:phyx-mini-0920:phyxminiproblems-problemphyxmini0920-answerchoice}
  The scalar expression printed beside each answer label, evaluated on coherent SI readouts. Choice B is transcribed faithfully even though its printed expression has the wrong physical dimension for an EMF
\end{definition}
\begin{proof}
  This is the declared model definition of displayed Expression In SI.
\end{proof}
\begin{definition}[recorded Dataset Answer]
  \label{def:physics:phyx-mini-0920:phyxminiproblems-problemphyxmini0920-recordeddatasetanswer}
  \lean{PhyXMiniProblems.ProblemPhyXMini0920.recordedDatasetAnswer}
  \uses{def:physics:phyx-mini-0920:phyxminiproblems-problemphyxmini0920-answerchoice}
  The answer label recorded in the supplied dataset metadata
\end{definition}
\begin{proof}
  This is the declared model definition of recorded Dataset Answer.
\end{proof}
\begin{lemma}[motional Emf Density eq angular Speed mul field mul radius]
  \label{lem:physics:phyx-mini-0920:phyxminiproblems-problemphyxmini0920-motionalemfdensity-eq-angularspeed-mul-field-mul-radius}
  \lean{PhyXMiniProblems.ProblemPhyXMini0920.motionalEmfDensity_eq_angularSpeed_mul_field_mul_radius}
  \uses{def:physics:phyx-mini-0920:phyxminiproblems-problemphyxmini0920-lengthinmeters, def:physics:phyx-mini-0920:phyxminiproblems-problemphyxmini0920-angularspeedinradianspersecond, def:physics:phyx-mini-0920:phyxminiproblems-problemphyxmini0920-magneticfluxdensityinteslas, def:physics:phyx-mini-0920:phyxminiproblems-problemphyxmini0920-motionalemfdensityinvoltspermeter, def:physics:phyx-mini-0920:phyxminiproblems-problemphyxmini0920-rotatingconductingdisksetup, def:physics:phyx-mini-0920:phyxminiproblems-problemphyxmini0920-satisfiesrotatingdiskmotionalemflaws}
  Combining the two local laws gives the differential coefficient dℰ/dr = ω B r on the disk. This is an intermediate consequence, not a premise containing the requested integrated answer
\end{lemma}
\begin{proof}
  The conclusion follows from the governing relations and figure readouts named in the statement of motional Emf Density eq angular Speed mul field mul radius.
\end{proof}
% --- Archon declaration topology end ---
% --- Archon physics formalization source end ---
